% ====================================================================================
\chapter*{附录C：大模型工程速成——显存、算力与实现}
\addcontentsline{toc}{chapter}{附录C：大模型工程速成——显存、算力与实现}
% ====================================================================================
% ====================================================================================

\section*{引言：从论文到实验室}

\sidenoteplain{本附录内容是你在实验室跑实验必须具备的"工程直觉"。不懂这些，代码写对了也跑不起来。}

在前面的章节中，我们深入理解了Transformer的数学原理。但当你真正要在实验室训练或部署模型时，会遇到一系列非常实际的问题：

\begin{itemize}
    \item "导师给了我一个7B的模型，实验室的显卡能跑吗？"
    \item "为什么我的训练刚开始就报OOM（Out of Memory）？"
    \item "同学说他用4张卡，我只有2张，我的实验要跑多久？"
\end{itemize}

这些问题的答案需要你具备\textbf{工程直觉}——能够快速估算显存、预测训练时间、写出正确的代码。
\clarify{本附录不引入新数学。它回答的是第1章那个问题在工程里的版本：显存、算力、带宽是硬约束，"影子价格"告诉你多买一张卡、多要一小时机时能换来多少。}

\begin{quote}
\textbf{本附录目标}：把"实验能不能跑、要跑多久"变成一个你会算的约束优化问题。目标还是损失函数，但约束换成了三种硬资源——显存、算力、带宽——它们的拉格朗日乘子就是"显存/算力/带宽的影子价格"：多买 1 GB 显存、多要 1 小时机时，损失能降多少。读完本附录，你能在一张纸上估出任何模型的显存、训练时间和推理吞吐量，并且明白梯度检查点、LoRA、KV Cache、FlashAttention、MoE 都是在这三种资源之间按影子价格做交换。
\end{quote}

本附录分为六个部分：
\begin{enumerate}
    \item \textbf{硬件基础}：了解你手里的GPU能做什么
    \item \textbf{显存估算}：判断实验能不能跑起来
    \item \textbf{吞吐量计算}：预估实验要跑多久
    \item \textbf{代码实现}：手写核心组件
    \item \textbf{高效微调}：LoRA——用秩约束换显存
    \item \textbf{高效计算}：FlashAttention 与 MoE——在算力、显存、带宽之间做交换
\end{enumerate}

\begin{tcolorbox}[colback=green!5, colframe=green!60!black, title=学完这章，你将能够...]
\begin{itemize}
    \item 估算任意模型的显存占用（参数量→字节数）
    \item 判断你的GPU能否跑某个大模型
    \item 估算训练时间：算力×数据量→小时数
    \item 用PyTorch从零实现Attention、FFN等核心模块
    \item 使用混合精度训练、梯度检查点等工程技巧
\end{itemize}
\end{tcolorbox}

\begin{tcolorbox}[colback=yellow!10, colframe=orange!80, title=本附录的核心观点]
\textbf{工程问题不在主线之外——它就是第1章的拉格朗日乘子法，只不过约束从"物理方程"变成了"显存 $\le$ 显卡容量、时间 $\le$ 机时预算、数据搬运 $\le$ 带宽"。}

本附录的问题就是第1章的问题，只是目标（损失）不变，约束换成了资源不等式；于是乘子变成资源的影子价格，"紧约束 / 松约束"决定了你该优化显存、算力还是带宽。

\begin{center}
\renewcommand{\arraystretch}{1.3}
\small
\begin{tabular}{lp{3.6cm}p{7cm}}
\toprule
 & \textbf{第1章} & \textbf{本附录} \\
\midrule
变量 & $x \in \R^n$ & 参数 $\theta$，以及 batch $B$、上下文 $T$、精度、LoRA 秩 $r$ \\
目标 & $f(x)$ & 训练 / 验证损失 \\
约束 & $h_i(x) \le 0$ & $M(B, T, N, \text{精度}) \le M_{\text{GPU}}$；$\dfrac{6ND}{\#\text{GPU} \times \text{算力} \times \text{MFU}} \le T_{\text{预算}}$；IO $\le$ 带宽 \\
乘子 & $\mu_i \ge 0$ & $\lambda_{\text{mem}}, \lambda_{\text{time}}, \lambda_{\text{bw}}$：三种资源的影子价格 \\
一阶条件 & KKT：$\mu_i h_i = 0$ & 紧约束 $\lambda > 0$（OOM、带宽受限的 decode）；松约束 $\lambda = 0$（MFU 只有 40\%） \\
"技巧" & 沿约束面移动 & 检查点 / FlashAttention（算力换显存）、KV Cache（显存换算力）、LoRA（秩约束换显存）、MoE（参数换算力） \\
\bottomrule
\end{tabular}
\end{center}
\end{tcolorbox}

% ====================================================================================
\section*{认识你的GPU}
\addcontentsline{toc}{section}{认识你的GPU}
% ====================================================================================

\storynote{黄仁勋}{1993年创立NVIDIA。2006年推出CUDA时，华尔街嘲笑他"给游戏卡加科学计算功能"；2012年AlexNet用两张GTX 580训练出来，深度学习浪潮由此开始。今天每一次OOM报错，都在提醒你他的显卡卖得有多贵。}

% ------------------------------------------------------------------------------------
\subsection*{常见GPU参数速查}
% ------------------------------------------------------------------------------------

作为学生，你可能接触到以下几类GPU：

\begin{table}[htbp]
\centering
\small
\renewcommand{\arraystretch}{1.3}
\begin{tabular}{l|c|c|c|c|c}
\toprule
\textbf{型号} & \textbf{显存} & \textbf{FP16算力} & \textbf{显存带宽} & \textbf{功耗} & \textbf{定位} \\
\midrule
\multicolumn{6}{c}{\textit{消费级（个人/小实验室常见）}} \\
\midrule
RTX 3080 & 10 GB & 30 TFLOPS & 760 GB/s & 320W & 入门深度学习 \\
RTX 3090 & 24 GB & 36 TFLOPS & 936 GB/s & 350W & 个人炼丹主力 \\
RTX 4080 & 16 GB & 49 TFLOPS & 717 GB/s & 320W & 性价比之选 \\
RTX 4090 & 24 GB & 83 TFLOPS & 1008 GB/s & 450W & 消费级天花板 \\
\midrule
\multicolumn{6}{c}{\textit{专业级（实验室/计算中心常见）}} \\
\midrule
A100 40GB & 40 GB & 312 TFLOPS & 1555 GB/s & 400W & 主流科研卡 \\
A100 80GB & 80 GB & 312 TFLOPS & 2039 GB/s & 400W & 大模型训练 \\
H100 SXM & 80 GB & 990 TFLOPS & 3350 GB/s & 700W & 最新旗舰 \\
\midrule
\multicolumn{6}{c}{\textit{国产替代（部分高校可能接触）}} \\
\midrule
华为910B & 64 GB & 320 TFLOPS & 1200 GB/s & 400W & 国产主力 \\
\bottomrule
\end{tabular}
\caption{常见GPU参数对比（FP16/BF16算力，2024年数据）}
\end{table}

\sidenoteplain{\textbf{TFLOPS}：每秒万亿次浮点运算。1 TFLOPS = $10^{12}$ FLOPS。这是衡量GPU计算能力的核心指标。}

\paragraph{几个关键概念}

\begin{itemize}
    \item \textbf{显存（VRAM）}：GPU的"内存"，决定能装多大的模型
    \item \textbf{算力（FLOPS）}：计算速度，决定训练/推理多快
    \item \textbf{显存带宽}：数据搬运速度，大模型推理的瓶颈往往在这里
    \item \textbf{功耗}：电费和散热的考量，实验室可能有限制
\end{itemize}

% ------------------------------------------------------------------------------------
\subsection*{数据精度与存储}
% ------------------------------------------------------------------------------------

不同的数值精度决定了模型占用的显存大小：

\begin{table}[htbp]
\centering
\renewcommand{\arraystretch}{1.3}
\begin{tabular}{l|c|c|c|l}
\toprule
\textbf{精度} & \textbf{位数} & \textbf{字节} & \textbf{数值范围} & \textbf{典型用途} \\
\midrule
FP32 & 32 bit & 4 B & $\pm 3.4 \times 10^{38}$ & 传统训练、优化器状态 \\
TF32 & 19 bit & 4 B & 同FP32 & A100默认训练精度 \\
FP16 & 16 bit & 2 B & $\pm 6.5 \times 10^{4}$ & 混合精度训练 \\
BF16 & 16 bit & 2 B & $\pm 3.4 \times 10^{38}$ & 大模型训练（推荐） \\
INT8 & 8 bit & 1 B & $-128 \sim 127$ & 量化推理 \\
INT4 & 4 bit & 0.5 B & $-8 \sim 7$ & 激进量化 \\
\bottomrule
\end{tabular}
\caption{常见数值精度对比}
\end{table}

\begin{figure}[htbp]
\centering
\begin{tikzpicture}[scale=0.9]
    % FP32
    \node[font=\bfseries] at (-2, 2) {FP32:};
    \draw (0, 1.7) rectangle (0.8, 2.3);
    \node[font=\scriptsize] at (0.4, 2) {符号};
    \draw (0.8, 1.7) rectangle (3.2, 2.3);
    \node[font=\scriptsize] at (2, 2) {指数 (8位)};
    \draw (3.2, 1.7) rectangle (8, 2.3);
    \node[font=\scriptsize] at (5.6, 2) {尾数 (23位)};
    
    % FP16
    \node[font=\bfseries] at (-2, 1) {FP16:};
    \draw (0, 0.7) rectangle (0.5, 1.3);
    \node[font=\scriptsize] at (0.25, 1) {符};
    \draw (0.5, 0.7) rectangle (2, 1.3);
    \node[font=\scriptsize] at (1.25, 1) {指数 (5位)};
    \draw (2, 0.7) rectangle (5, 1.3);
    \node[font=\scriptsize] at (3.5, 1) {尾数 (10位)};
    
    % BF16
    \node[font=\bfseries] at (-2, 0) {BF16:};
    \draw (0, -0.3) rectangle (0.5, 0.3);
    \node[font=\scriptsize] at (0.25, 0) {符};
    \draw (0.5, -0.3) rectangle (2.9, 0.3);
    \node[font=\scriptsize] at (1.7, 0) {指数 (8位)};
    \draw (2.9, -0.3) rectangle (5, 0.3);
    \node[font=\scriptsize] at (3.95, 0) {尾数 (7位)};
    
    % 说明
    \node[right, font=\small, gray] at (8.5, 2) {精度最高，但慢且费显存};
    \node[right, font=\small, gray] at (5.5, 1) {容易溢出（范围小）};
    \node[right, font=\small, gray] at (5.5, 0) {范围大，精度够用（推荐）};
\end{tikzpicture}
\caption{不同浮点格式的位分配}
\end{figure}

\begin{keyidea}[BF16 vs FP16]
\begin{itemize}
    \item \textbf{FP16}：尾数多（精度高），但指数少（范围小），训练时容易溢出
    \item \textbf{BF16}：指数与FP32相同（范围大），尾数少但够用，训练更稳定
    \item \textbf{实践建议}：如果你的GPU支持BF16（A100、H100、RTX 30/40系列），优先用BF16
\end{itemize}
\end{keyidea}

% ====================================================================================
\section*{显存估算：你的实验能跑吗？}
\addcontentsline{toc}{section}{显存估算：你的实验能跑吗？}
% ====================================================================================

显存是深度学习的"硬通货"。显存不够，代码写得再对也跑不起来。
\review{第1章}{"显存是硬通货"就是说显存是紧约束：$\lambda_{\text{mem}} > 0$。第1章讲过，紧约束的乘子等于"放松一单位约束，目标能改善多少"——这里就是"多 1 GB 显存能换来多大的 batch 或上下文"。}

% ------------------------------------------------------------------------------------
\subsection*{模型参数的显存占用}
% ------------------------------------------------------------------------------------

\begin{keyformula}[title=模型显存速算]
\begin{equation}
    \text{模型显存 (GB)} = \text{参数量 (B)} \times \text{每参数字节数}
\end{equation}

速算口诀：
\begin{itemize}
    \item FP32：参数量 $\times$ 4
    \item FP16/BF16：参数量 $\times$ 2
    \item INT8：参数量 $\times$ 1
    \item INT4：参数量 $\times$ 0.5
\end{itemize}
\end{keyformula}

\begin{table}[htbp]
\centering
\renewcommand{\arraystretch}{1.3}
\begin{tabular}{l|c|c|c|c|c}
\toprule
\textbf{模型} & \textbf{参数量} & \textbf{FP32} & \textbf{FP16/BF16} & \textbf{INT8} & \textbf{INT4} \\
\midrule
GPT-2 & 1.5B & 6 GB & 3 GB & 1.5 GB & 0.75 GB \\
LLaMA-7B & 7B & 28 GB & 14 GB & 7 GB & 3.5 GB \\
LLaMA-13B & 13B & 52 GB & 26 GB & 13 GB & 6.5 GB \\
Qwen-14B & 14B & 56 GB & 28 GB & 14 GB & 7 GB \\
LLaMA-70B & 70B & 280 GB & 140 GB & 70 GB & 35 GB \\
\bottomrule
\end{tabular}
\caption{常见模型的参数显存需求（仅模型权重）}
\end{table}

\begin{example}[判断能否跑LLaMA-7B推理]
\textbf{场景}：实验室有一张RTX 3090（24GB），想跑LLaMA-7B推理。

\textbf{分析}：
\begin{itemize}
    \item FP16加载：$7 \times 2 = 14$ GB $\checkmark$ 可以
    \item 还剩 $24 - 14 = 10$ GB 给KV Cache和其他开销
    \item 结论：\textbf{可以跑}，但长文本生成时KV Cache可能爆显存
\end{itemize}

\textbf{如果只有RTX 3080（10GB）}：
\begin{itemize}
    \item FP16加载：14 GB > 10 GB $\times$ 跑不了
    \item INT4量化：$7 \times 0.5 = 3.5$ GB $\checkmark$ 可以
    \item 结论：必须用\textbf{4-bit量化}（如GPTQ、AWQ）
\end{itemize}
\end{example}

% ------------------------------------------------------------------------------------
\subsection*{训练显存：为什么比推理大这么多？}
% ------------------------------------------------------------------------------------

新手常见困惑：为什么推理只要14GB，训练却要上百GB？

答案：训练需要存储的东西远不止模型权重。

\begin{figure}[htbp]
\centering
\begin{tikzpicture}[>=stealth, scale=0.85]
    % 推理
    \begin{scope}[xshift=-4.5cm]
        \node[font=\bfseries] at (1, 5.5) {推理};
        
        \draw[fill=blue!30] (0, 0) rectangle (2, 3);
        \node at (1, 1.5) {模型权重};
        \node[font=\small] at (1, 0.8) {14 GB};
        
        \draw[fill=gray!20] (0, 3) rectangle (2, 4);
        \node[font=\small] at (1, 3.5) {KV Cache};
        
        \draw[<->] (2.5, 0) -- (2.5, 4);
        \node[right, font=\small] at (2.5, 2) {$\approx$16 GB};
    \end{scope}
    
    % 训练
    \begin{scope}[xshift=3cm]
        \node[font=\bfseries] at (1.5, 5.5) {全量微调};
        
        \draw[fill=blue!30] (0, 0) rectangle (3, 1.5);
        \node at (1.5, 0.75) {模型权重 (FP16)};
        \node[font=\small, gray] at (1.5, 0.3) {14 GB};
        
        \draw[fill=green!30] (0, 1.5) rectangle (3, 3);
        \node at (1.5, 2.25) {梯度 (FP16)};
        \node[font=\small, gray] at (1.5, 1.8) {14 GB};
        
        \draw[fill=orange!30] (0, 3) rectangle (3, 6);
        \node at (1.5, 4.5) {优化器状态};
        \node[font=\small] at (1.5, 4) {(Adam, FP32)};
        \node[font=\small, gray] at (1.5, 3.5) {56 GB};
        
        \draw[fill=red!20] (0, 6) rectangle (3, 7.5);
        \node at (1.5, 6.75) {激活值};
        \node[font=\small, gray] at (1.5, 6.3) {可变};
        
        \draw[<->] (3.5, 0) -- (3.5, 6);
        \node[right, font=\small] at (3.5, 3) {$\geq$84 GB};
        \node[right, font=\scriptsize, red] at (3.5, 2.3) {一张A100都不够！};
    \end{scope}
\end{tikzpicture}
\caption{LLaMA-7B推理 vs 训练的显存对比}
\end{figure}

\paragraph{训练显存的四个组成部分}

\begin{enumerate}
    \item \textbf{模型权重}：和推理一样
    \begin{equation}
        M_{\text{weights}} = N \times \text{dtype\_size}
    \end{equation}
    
    \item \textbf{梯度}：每个参数对应一个梯度，大小与权重相同
    \begin{equation}
        M_{\text{gradients}} = N \times \text{dtype\_size}
    \end{equation}
    
    \item \textbf{优化器状态}：这是大头！以Adam为例：
    \begin{itemize}
        \item 一阶矩（Momentum）：$m_t = \beta_1 m_{t-1} + (1-\beta_1) g_t$
        \item 二阶矩（Variance）：$v_t = \beta_2 v_{t-1} + (1-\beta_2) g_t^2$
        \item 通常用FP32存储以保证数值稳定
    \end{itemize}
    \begin{equation}
        M_{\text{optimizer}} = N \times 2 \times 4 = 8N \quad \text{(Adam)}
    \end{equation}
    
    \item \textbf{激活值}：前向传播的中间结果，反向时要用
    \begin{itemize}
        \item 取决于Batch Size、序列长度、模型深度
        \item 可以用\textbf{梯度检查点}节省（用计算换显存）
        \clarify{梯度检查点（\texttt{torch.utils.checkpoint}）：前向只保存每隔几层的激活值，反向时重新计算中间层。激活显存从 $O(L)$ 降到 $O(\sqrt{L})$，代价是多做约三分之一的前向——这是在用算力的影子价格换显存的影子价格。}
    \end{itemize}
\end{enumerate}

\begin{keyformula}[title=全量微调显存估算]
\begin{equation}
    M_{\text{train}} \approx \underbrace{2N}_{\text{权重}} + \underbrace{2N}_{\text{梯度}} + \underbrace{8N}_{\text{优化器}} + M_{\text{激活}} = 12N + M_{\text{激活}} \quad \text{（混合精度若另存 FP32 主权重，再加 } 4N\text{）}
\end{equation}

简化记忆：\textbf{训练显存 $\approx$ 推理显存 $\times$ 6倍}（不含激活值）
\end{keyformula}

\begin{table}[htbp]
\centering
\renewcommand{\arraystretch}{1.3}
\begin{tabular}{l|c|c|c|c}
\toprule
\textbf{模型} & \textbf{FP16推理} & \textbf{全量微调} & \textbf{单卡可行性} & \textbf{解决方案} \\
\midrule
GPT-2 (1.5B) & 3 GB & $\approx$20 GB & RTX 3090 \checkmark{} & 单卡可跑 \\
LLaMA-7B & 14 GB & $\approx$100 GB & A100-80G $\times$ & DeepSpeed / LoRA \\
LLaMA-13B & 26 GB & $\approx$180 GB & 单卡不可能 & 多卡并行 \\
LLaMA-70B & 140 GB & $\approx$1 TB & 单卡不可能 & 集群训练 \\
\bottomrule
\end{tabular}
\caption{不同模型的训练显存需求与可行性}
\end{table}

% ------------------------------------------------------------------------------------
\subsection*{KV Cache：长文本的显存杀手}
% ------------------------------------------------------------------------------------

自回归生成时，为避免重复计算，需要缓存历史token的Key和Value。

\begin{keyformula}[title=KV Cache显存]
\begin{equation}
    M_{KV} = 2 \times L \times S \times H \times D_{\text{type}}
\end{equation}
\begin{itemize}
    \item 2：K和V各一份
    \item $L$：Transformer层数
    \item $S$：序列长度（输入 + 已生成）
    \item $H$：隐藏层维度
    \item $D_{\text{type}}$：每元素字节数（FP16为2）
\end{itemize}
\end{keyformula}

\begin{example}[KV Cache计算]
\textbf{配置}：LLaMA-7B（32层，$H=4096$），序列长度4K，Batch=1（即同时处理1个请求/1个并发用户），FP16

\begin{align}
    M_{KV} &= 2 \times 32 \times 4096 \times 4096 \times 2 \text{ Bytes} \\
    &= 2,147,483,648 \text{ Bytes} \approx 2 \text{ GB}
\end{align}

\textbf{追问}：如果Batch Size变成16呢？
\begin{equation}
    M_{KV} = 2 \text{ GB} \times 16 = 32 \text{ GB}
\end{equation}

这就是为什么大Batch + 长序列推理这么吃显存。
\end{example}
\think{把 KV Cache 和下一节的 FLOPs 公式合起来看：缓存省的是算力，花的是显存。什么时候不该缓存？（提示：batch 很大、上下文很短、显存已经是紧约束时。）}

\begin{table}[htbp]
\centering
\small
\renewcommand{\arraystretch}{1.3}
\begin{tabular}{l|c|c|c|c}
\toprule
\textbf{序列长度} & \textbf{Batch=1} & \textbf{Batch=8} & \textbf{Batch=32} & \textbf{Batch=64} \\
\midrule
2K & 1 GB & 8 GB & 32 GB & 64 GB \\
4K & 2 GB & 16 GB & 64 GB & 128 GB \\
8K & 4 GB & 32 GB & 128 GB & 256 GB \\
32K & 16 GB & 128 GB & 512 GB & 1 TB \\
128K & 64 GB & 512 GB & 2 TB & 4 TB \\
\bottomrule
\end{tabular}
\caption{LLaMA-7B的KV Cache显存（FP16）}
\end{table}

\paragraph{缓解KV Cache的方法}

\begin{itemize}
    \item \textbf{KV Cache量化}：用INT8存储K和V，显存减半
    \item \textbf{GQA}（Grouped Query Attention）：多个Query头共享K/V，LLaMA-2/3已采用
    \item \textbf{PagedAttention}：动态分配KV Cache，vLLM使用的技术
    \item \textbf{滑动窗口}：只保留最近的K/V，Mistral采用
\end{itemize}

% ------------------------------------------------------------------------------------
\subsection*{显存速查表}
% ------------------------------------------------------------------------------------

\begin{table}[htbp]
\centering
\small
\renewcommand{\arraystretch}{1.4}
\begin{tabular}{l|l|l}
\toprule
\textbf{场景} & \textbf{速算公式} & \textbf{7B模型示例} \\
\midrule
FP16推理 & $2 \times N$ & $2 \times 7 = 14$ GB \\
FP32推理 & $4 \times N$ & $4 \times 7 = 28$ GB \\
INT8推理 & $1 \times N$ & $1 \times 7 = 7$ GB \\
INT4推理 & $0.5 \times N$ & $0.5 \times 7 = 3.5$ GB \\
\midrule
全量微调 & $\approx 12 \times N$ & $12 \times 7 = 84$ GB \\
LoRA微调 & $\approx 2 \times N$ & $\approx 16$ GB \\
\midrule
KV Cache (FP16) & $4 \times L \times S \times H$ & 见上表 \\
\bottomrule
\end{tabular}
\caption{显存估算速查表（$N$单位：B参数，结果单位：GB）}
\end{table}

% ====================================================================================
\section*{吞吐量计算：实验要跑多久？}
\addcontentsline{toc}{section}{吞吐量计算：实验要跑多久？}
% ====================================================================================

显存够了之后，下一个问题是：\textbf{实验要跑多久？}

% ------------------------------------------------------------------------------------
\subsection*{计算量估算}
% ------------------------------------------------------------------------------------

\paragraph{Transformer的计算量}

对于一个Transformer模型，前向传播的计算量约为：

\begin{keyformula}[title=Transformer计算量估算]
\begin{equation}
    \text{FLOPs}_{\text{forward}} \approx 2 \times N \quad \text{（每个 token）}
\end{equation}
其中 $N$ 是参数量；处理 $S$ 个 token 就是 $2NS$。

训练时还需要反向传播（约等于2倍前向），所以：
\begin{equation}
    \text{FLOPs}_{\text{train}} \approx 6 \times N \ \text{（每个 token）}, \qquad \text{训练 } D \text{ 个 token 共 } 6ND
\end{equation}
\end{keyformula}
\review{第8章}{为什么反向约等于 2 倍前向？每个 $y = Wx$ 反向要算两个乘法：$\partial L/\partial x = W^\top \lambda$（伴随方程）和 $\partial L/\partial W = \lambda x^\top$（参数梯度）。第8章的伴随方程正是这两项。}

\sidenoteplain{这个公式假设模型以矩阵乘法为主，忽略了激活函数等开销。实际计算量可能略高。}

\begin{example}[训练计算量估算]
\textbf{任务}：用100B tokens预训练一个7B模型

\textbf{计算}：
\begin{align}
    \text{总FLOPs} &= 6 \times 7 \times 10^9 \times 100 \times 10^9 \\
    &= 4.2 \times 10^{21} \text{ FLOPs} \\
    &= 4.2 \text{ ZFLOPs}
\end{align}
\end{example}

% ------------------------------------------------------------------------------------
\subsection*{训练时间估算}
% ------------------------------------------------------------------------------------

知道计算量后，可以估算训练时间：

\begin{keyformula}[title=训练时间估算]
\begin{equation}
    T = \frac{\text{总FLOPs}}{\text{GPU数量} \times \text{单卡算力} \times \text{利用率}}
\end{equation}
\end{keyformula}

\paragraph{GPU利用率（MFU）}

实际训练中，GPU不可能100\%满载。\textbf{模型算力利用率}（Model FLOPS Utilization, MFU）通常在30\%-60\%：

\begin{table}[htbp]
\centering
\renewcommand{\arraystretch}{1.3}
\begin{tabular}{l|c|l}
\toprule
\textbf{场景} & \textbf{典型MFU} & \textbf{原因} \\
\midrule
单卡训练小模型 & 40-50\% & 数据加载、小batch效率低 \\
单卡训练大模型 & 50-60\% & batch大，计算密集 \\
多卡数据并行 & 40-55\% & 通信开销 \\
多卡模型并行 & 30-45\% & 流水线气泡、通信开销 \\
\bottomrule
\end{tabular}
\caption{不同场景的典型GPU利用率}
\end{table}

\begin{example}[预训练时间估算]
\textbf{任务}：用100B tokens预训练7B模型

\textbf{配置}：8张A100-80GB，MFU假设50\%

\textbf{计算}：
\begin{align}
    \text{有效算力} &= 8 \times 312 \text{ TFLOPS} \times 50\% = 1248 \text{ TFLOPS} \\
    T &= \frac{4.2 \times 10^{21}}{1248 \times 10^{12}} = 3.37 \times 10^{6} \text{ 秒} \\
    &\approx 39 \text{ 天}
\end{align}

\textbf{如果用64张H100}：
\begin{align}
    \text{有效算力} &= 64 \times 990 \text{ TFLOPS} \times 50\% = 31680 \text{ TFLOPS} \\
    T &= \frac{4.2 \times 10^{21}}{31680 \times 10^{12}} \approx 1.5 \text{ 天}
\end{align}
\end{example}

% ------------------------------------------------------------------------------------
\subsection*{不同规模的训练成本参考}
% ------------------------------------------------------------------------------------

\begin{table}[htbp]
\centering
\small
\renewcommand{\arraystretch}{1.4}
\begin{tabular}{l|c|c|c|c}
\toprule
\textbf{模型} & \textbf{训练数据} & \textbf{总FLOPs} & \textbf{8×A100时间} & \textbf{64×H100时间} \\
\midrule
GPT-2 (1.5B) & 40B tokens & $3.6 \times 10^{20}$ & 3天 & 4小时 \\
LLaMA-7B & 1T tokens & $4.2 \times 10^{22}$ & 1年+ & 15天 \\
LLaMA-13B & 1T tokens & $7.8 \times 10^{22}$ & 2年+ & 28天 \\
LLaMA-70B & 1.4T tokens & $5.9 \times 10^{23}$ & 15年+ & 7个月 \\
\bottomrule
\end{tabular}
\caption{预训练时间估算（MFU=50\%）}
\end{table}

\begin{remark}
从这个表可以看出：
\begin{itemize}
    \item 预训练大模型需要\textbf{大量算力}，不是学生个人能承担的
    \item 学生能做的是\textbf{微调}（只需要原训练的千分之一算力）
    \item 或者\textbf{小模型实验}（验证想法后再申请资源）
\end{itemize}
\end{remark}
\preview{第14章 PINN 和第15章神经算子的实验都在"单卡几小时"量级——先用本节公式算一遍，再决定要不要申请集群。}

% ------------------------------------------------------------------------------------
\subsection*{微调时间估算}
% ------------------------------------------------------------------------------------

微调比预训练快得多，因为数据量小：

\begin{example}[LoRA微调时间估算]
\textbf{任务}：对LLaMA-7B做LoRA微调，10K条指令数据，平均长度512 tokens

\textbf{计算}：
\begin{align}
    \text{总tokens} &= 10000 \times 512 = 5.12 \times 10^6 \\
    \text{每epoch FLOPs} &= 6 \times 7 \times 10^9 \times 5.12 \times 10^6 = 2.15 \times 10^{17}
\end{align}

\textbf{单张RTX 4090}（83 TFLOPS，MFU 45\%）：
\begin{align}
    \text{每epoch时间} &= \frac{2.15 \times 10^{17}}{83 \times 10^{12} \times 0.45} \approx 5760 \text{ 秒} \approx 1.6 \text{ 小时}
\end{align}

训练3个epoch：\textbf{约5小时}
\end{example}

% ------------------------------------------------------------------------------------
\subsection*{推理吞吐量}
% ------------------------------------------------------------------------------------

推理时，我们关心的是\textbf{每秒能生成多少token}（tokens/s）。

\paragraph{推理的两个阶段}

\begin{enumerate}
    \item \textbf{Prefill}（预填充）：处理输入的所有token，计算密集
    \item \textbf{Decode}（解码）：逐个生成输出token，显存带宽受限
\end{enumerate}

\begin{figure}[htbp]
\centering
\begin{tikzpicture}[>=stealth, scale=0.85]
    % Prefill
    \draw[fill=blue!20] (0, 0) rectangle (3, 2);
    \node at (1.5, 1) {Prefill};
    \node[below, font=\small] at (1.5, 0) {处理输入};
    
    % Decode
    \draw[fill=orange!20] (4, 0) rectangle (10, 2);
    \node at (7, 1) {Decode（自回归生成）};
    \node[below, font=\small] at (7, 0) {逐token生成};
    
    % 时间轴
    \draw[->] (0, -1) -- (11, -1) node[right] {时间};
    
    % 特点标注
    \node[above, font=\small, blue] at (1.5, 2.2) {计算密集};
    \node[above, font=\small, blue] at (1.5, 2.7) {可高度并行};
    
    \node[above, font=\small, orange] at (7, 2.2) {带宽受限};
    \node[above, font=\small, orange] at (7, 2.7) {逐个串行};
\end{tikzpicture}
\caption{推理的两个阶段}
\end{figure}

\paragraph{Decode阶段的吞吐量估算}

Decode阶段是带宽受限的，每生成一个token需要从显存读取整个模型：

\begin{equation}
    \text{最大tokens/s} \approx \frac{\text{显存带宽}}{\text{模型大小}}
\end{equation}

\begin{example}[LLaMA-7B推理吞吐量]
\textbf{配置}：单张A100-80GB（带宽2039 GB/s），FP16

\begin{align}
    \text{理论最大} &= \frac{2039 \text{ GB/s}}{14 \text{ GB}} \approx 145 \text{ tokens/s}
\end{align}

实际由于各种开销，通常能达到\textbf{80-100 tokens/s}。

\textbf{对比RTX 4090}（带宽1008 GB/s）：
\begin{align}
    \text{理论最大} &= \frac{1008}{14} \approx 72 \text{ tokens/s}
\end{align}
实际约\textbf{40-50 tokens/s}。
\end{example}

\begin{table}[htbp]
\centering
\small
\renewcommand{\arraystretch}{1.3}
\begin{tabular}{l|c|c|c|c}
\toprule
\textbf{GPU} & \textbf{LLaMA-7B} & \textbf{LLaMA-13B} & \textbf{LLaMA-70B} & \textbf{备注} \\
\midrule
RTX 3090 & 30-40 t/s & 需量化 & 不可行 & 消费级 \\
RTX 4090 & 40-50 t/s & 20-25 t/s & 需量化 & 消费级天花板 \\
A100-40GB & 70-90 t/s & 40-50 t/s & 需多卡 & 专业卡 \\
A100-80GB & 80-100 t/s & 50-60 t/s & 15-20 t/s & 大显存 \\
H100 & 150-200 t/s & 100-120 t/s & 40-50 t/s & 最新旗舰 \\
\bottomrule
\end{tabular}
\caption{不同GPU的LLM推理吞吐量参考（FP16，Batch=1，实测范围）}
\end{table}

% ====================================================================================
\section*{代码实现：手写Attention}
\addcontentsline{toc}{section}{代码实现：手写Attention}
% ====================================================================================

\sidenoteplain{\textbf{Transformer的诞生}：2017年，Google的Vaswani等人发表``Attention Is All You Need''，用纯注意力机制取代了RNN/CNN，成为现代大模型的基石。论文标题成为深度学习史上最经典的口号之一。}

理解原理后，让我们动手实现核心组件。

% ------------------------------------------------------------------------------------
\subsection*{Scaled Dot-Product Attention}
% ------------------------------------------------------------------------------------

\sidenoteplain{\textbf{注意力机制溯源}：注意力机制最早由Bahdanau等人(2014)引入机器翻译，让模型``看向''输入的不同部分。Transformer将其推广为``自注意力''，让序列中的每个位置都能关注其他所有位置。}

回顾公式：
\begin{equation}
    \text{Attention}(Q, K, V) = \text{softmax}\left(\frac{QK^\top}{\sqrt{d_k}}\right) V
\end{equation}

\begin{lstlisting}[language=Python, caption=Scaled Dot-Product Attention]
import torch
import torch.nn as nn
import torch.nn.functional as F
import math

def scaled_dot_product_attention(query, key, value, mask=None):
    """
    缩放点积注意力
    
    Args:
        query: [batch, num_heads, seq_len_q, head_dim]
        key:   [batch, num_heads, seq_len_k, head_dim]
        value: [batch, num_heads, seq_len_k, head_dim]
        mask:  [batch, 1, seq_len_q, seq_len_k] 或可广播形状
               1表示保留，0表示遮挡
    
    Returns:
        output: [batch, num_heads, seq_len_q, head_dim]
        attn_weights: [batch, num_heads, seq_len_q, seq_len_k]
    """
    d_k = query.size(-1)  # head_dim
    
    # Step 1: Q @ K^T -> [batch, heads, seq_q, seq_k]
    scores = torch.matmul(query, key.transpose(-2, -1))
    
    # Step 2: 缩放，防止点积过大导致softmax饱和
    scores = scores / math.sqrt(d_k)
    
    # Step 3: 应用mask（可选）
    if mask is not None:
        # mask为0的位置填负无穷，softmax后变成0
        scores = scores.masked_fill(mask == 0, float('-inf'))
    
    # Step 4: softmax归一化（在最后一个维度，即key的位置）
    attn_weights = F.softmax(scores, dim=-1)
    
    # Step 5: 加权求和
    output = torch.matmul(attn_weights, value)
    
    return output, attn_weights
\end{lstlisting}

\begin{example}[运行示例：理解张量维度变化]
让我们用具体数字追踪一次Attention计算：

\begin{lstlisting}[language=Python, caption=Attention维度演示]
# 输入: batch=2, heads=4, seq_len=3, head_dim=8
query = torch.randn(2, 4, 3, 8)
key   = torch.randn(2, 4, 3, 8)
value = torch.randn(2, 4, 3, 8)

output, weights = scaled_dot_product_attention(query, key, value)

print(f"Query shape:   {query.shape}")   # [2, 4, 3, 8]
print(f"Key shape:     {key.shape}")     # [2, 4, 3, 8]
print(f"Scores shape:  {(query @ key.transpose(-2,-1)).shape}")  # [2, 4, 3, 3]
print(f"Weights shape: {weights.shape}") # [2, 4, 3, 3] <- softmax后的注意力权重
print(f"Output shape:  {output.shape}")  # [2, 4, 3, 8]

# 验证：每行注意力权重和为1
print(f"权重和: {weights[0,0,0,:].sum():.4f}")  # 1.0000
\end{lstlisting}

\textbf{矩阵运算追踪}：
\begin{enumerate}
    \item $Q \times K^\top$: $[2,4,3,\mathbf{8}] \times [2,4,\mathbf{8},3] = [2,4,3,3]$ —— 每对位置的相似度
    \item softmax: $[2,4,3,3] \to [2,4,3,3]$ —— 归一化为概率
    \item $\text{weights} \times V$: $[2,4,3,\mathbf{3}] \times [2,4,\mathbf{3},8] = [2,4,3,8]$ —— 加权求和
\end{enumerate}

\textbf{核心洞察}：\textbf{Attention就是三次矩阵乘法}。这正是GPU能高效处理的原因！
\end{example}

\paragraph{关键点解析}

\begin{enumerate}
    \item \textbf{为什么除以$\sqrt{d_k}$？}
    
    假设Q和K的元素独立同分布于$\mathcal{N}(0, 1)$，则点积的方差是$d_k$：
    \begin{equation}
        \text{Var}(q \cdot k) = \sum_{i=1}^{d_k} \text{Var}(q_i k_i) = d_k
    \end{equation}
    
    $d_k$越大，点积值越大，softmax越接近one-hot，梯度越小。缩放后方差为1。
    
    \item \textbf{为什么mask用负无穷？}
    
    因为$\text{softmax}(-\infty) = 0$。如果直接把scores设成0，softmax后不是0而是一个正数。
    
    \item \textbf{softmax在哪个维度？}
    
    \texttt{dim=-1}，即对所有Key位置归一化，保证注意力权重和为1。
\end{enumerate}

% ------------------------------------------------------------------------------------
\subsection*{因果掩码（Causal Mask）}
% ------------------------------------------------------------------------------------

自回归生成时，位置$i$只能看到位置$j \leq i$的token：

\begin{lstlisting}[language=Python, caption=生成因果掩码]
def create_causal_mask(seq_len, device='cpu'):
    """
    创建因果掩码（下三角矩阵）
    
    Returns:
        mask: [1, 1, seq_len, seq_len]
    """
    # torch.tril: 下三角矩阵
    mask = torch.tril(torch.ones(seq_len, seq_len, device=device))
    # 扩展维度以便广播到 [batch, heads, seq, seq]
    return mask.unsqueeze(0).unsqueeze(0)

# 示例: seq_len=4
# [[1, 0, 0, 0],
#  [1, 1, 0, 0],
#  [1, 1, 1, 0],
#  [1, 1, 1, 1]]
\end{lstlisting}

\begin{figure}[htbp]
\centering
\begin{tikzpicture}[scale=0.55]
    % 矩阵格子
    \foreach \i in {0,...,5} {
        \draw[gray] (0, \i) -- (6, \i);
        \draw[gray] (\i, 0) -- (\i, 6);
    }
    \draw[thick] (0, 0) rectangle (6, 6);
    
    % 填充
    \foreach \i in {0,...,5} {
        \foreach \j in {0,...,5} {
            \pgfmathtruncatemacro{\keep}{\j <= \i ? 1 : 0}
            \ifnum\keep=1
                \fill[green!30] (\j, 5-\i) rectangle (\j+1, 6-\i);
                \node[font=\small] at (\j+0.5, 5.5-\i) {1};
            \else
                \fill[red!20] (\j, 5-\i) rectangle (\j+1, 6-\i);
                \node[font=\small] at (\j+0.5, 5.5-\i) {0};
            \fi
        }
    }
    
    % 标签
    \node[above] at (3, 6.3) {Key位置 $j$};
    \node[left, rotate=90] at (-0.5, 3) {Query位置 $i$};
    
    % 说明
    \node[right, font=\small] at (7, 5) {绿色：可以看到};
    \node[right, font=\small] at (7, 4) {红色：不能看到};
    \node[right, font=\small] at (7, 2.5) {位置 $i$ 只能attend到};
    \node[right, font=\small] at (7, 1.8) {位置 $j \leq i$};
\end{tikzpicture}
\caption{因果掩码示例（seq\_len=6）}
\end{figure}

% ------------------------------------------------------------------------------------
\subsection*{Multi-Head Attention的维度变换}
% ------------------------------------------------------------------------------------

Multi-Head的关键是正确处理维度：

\begin{equation}
    [B, L, D] \xrightarrow{\text{拆分}} [B, H, L, D_h] \xrightarrow{\text{attention}} [B, H, L, D_h] \xrightarrow{\text{合并}} [B, L, D]
\end{equation}

其中$D = H \times D_h$（总维度 = 头数 $\times$ 每头维度）。

\begin{lstlisting}[language=Python, caption=Multi-Head维度变换]
def split_heads(x, num_heads):
    """
    拆分成多头
    [batch, seq_len, hidden_size] -> [batch, num_heads, seq_len, head_dim]
    """
    batch, seq_len, hidden = x.size()
    head_dim = hidden // num_heads
    
    # reshape: [B, L, H] -> [B, L, num_heads, head_dim]
    x = x.view(batch, seq_len, num_heads, head_dim)
    
    # transpose: [B, L, num_heads, head_dim] -> [B, num_heads, L, head_dim]
    # 把num_heads提前，方便并行计算
    return x.transpose(1, 2)


def merge_heads(x):
    """
    合并多头
    [batch, num_heads, seq_len, head_dim] -> [batch, seq_len, hidden_size]
    """
    batch, num_heads, seq_len, head_dim = x.size()
    
    # transpose回去: [B, num_heads, L, head_dim] -> [B, L, num_heads, head_dim]
    x = x.transpose(1, 2)
    
    # reshape: [B, L, num_heads, head_dim] -> [B, L, hidden]
    # contiguous()保证内存连续，否则view可能报错
    return x.contiguous().view(batch, seq_len, num_heads * head_dim)
\end{lstlisting}

\paragraph{为什么要transpose？}

\texttt{torch.matmul}对最后两个维度做矩阵乘法，前面的维度广播。

把\texttt{num\_heads}放到前面后，每个head的$Q \in \mathbb{R}^{L \times D_h}$和$K \in \mathbb{R}^{L \times D_h}$可以并行相乘。

% ------------------------------------------------------------------------------------
\subsection*{完整的Multi-Head Attention}
% ------------------------------------------------------------------------------------

\begin{lstlisting}[language=Python, caption=完整的Multi-Head Attention]
class MultiHeadAttention(nn.Module):
    def __init__(self, hidden_size, num_heads, dropout=0.1):
        super().__init__()
        assert hidden_size % num_heads == 0, "hidden_size必须能被num_heads整除"
        
        self.hidden_size = hidden_size
        self.num_heads = num_heads
        self.head_dim = hidden_size // num_heads
        
        # Q, K, V投影
        self.q_proj = nn.Linear(hidden_size, hidden_size)
        self.k_proj = nn.Linear(hidden_size, hidden_size)
        self.v_proj = nn.Linear(hidden_size, hidden_size)
        
        # 输出投影
        self.out_proj = nn.Linear(hidden_size, hidden_size)
        self.dropout = nn.Dropout(dropout)
        
    def forward(self, x, mask=None):
        """
        自注意力：Q, K, V都来自同一个输入x
        Args:
            x: [batch, seq_len, hidden_size]
            mask: [batch, 1, seq_len, seq_len] 或 None
        """
        batch_size, seq_len, _ = x.size()
        
        # 1. 投影
        Q = self.q_proj(x)  # [B, L, H]
        K = self.k_proj(x)
        V = self.v_proj(x)
        
        # 2. 拆分多头
        Q = self._split_heads(Q)  # [B, num_heads, L, head_dim]
        K = self._split_heads(K)
        V = self._split_heads(V)
        
        # 3. 计算注意力
        scores = torch.matmul(Q, K.transpose(-2, -1)) / math.sqrt(self.head_dim)
        
        if mask is not None:
            scores = scores.masked_fill(mask == 0, float('-inf'))
        
        attn_weights = F.softmax(scores, dim=-1)
        attn_weights = self.dropout(attn_weights)
        
        attn_output = torch.matmul(attn_weights, V)  # [B, num_heads, L, head_dim]
        
        # 4. 合并多头
        attn_output = self._merge_heads(attn_output)  # [B, L, H]
        
        # 5. 输出投影
        output = self.out_proj(attn_output)
        
        return output
    
    def _split_heads(self, x):
        B, L, _ = x.size()
        return x.view(B, L, self.num_heads, self.head_dim).transpose(1, 2)
    
    def _merge_heads(self, x):
        B, _, L, _ = x.size()
        return x.transpose(1, 2).contiguous().view(B, L, self.hidden_size)
\end{lstlisting}

% ====================================================================================
\section*{高效微调：LoRA简介}
\addcontentsline{toc}{section}{高效微调：LoRA简介}
% ====================================================================================

\sidenoteplain{\textbf{LoRA的诞生}：2021年，微软的Hu等人提出LoRA，让普通研究者也能在消费级显卡上微调大模型。这篇论文引用量已超过5000次，催生了QLoRA、DoRA等一系列变体，成为开源社区微调LLM的标准方法。}

全量微调大模型需要的显存太大。\textbf{LoRA}（Low-Rank Adaptation）是一种参数高效的替代方案。

% ------------------------------------------------------------------------------------
\subsection*{核心思想}
% ------------------------------------------------------------------------------------

\begin{keyidea}[LoRA的假设]
微调时，权重的更新量$\Delta W$是\textbf{低秩}的。

原来的大矩阵$W \in \mathbb{R}^{d \times d}$可能有1600万参数，但更新量可以用两个小矩阵近似：
\begin{equation}
    \Delta W \approx BA, \quad B \in \mathbb{R}^{d \times r}, \quad A \in \mathbb{R}^{r \times d}
\end{equation}
其中$r$是秩，通常$r = 4, 8, 16$，远小于$d$。
\end{keyidea}
\review{第1章}{LoRA 是给优化问题加一条约束 $\mathrm{rank}(\Delta W) \le r$。第1章说过加约束会让最优值变差，但这里它换来的是显存约束从 $12N$ 降到约 $2N$——两条约束的影子价格之比决定 $r$ 该取多大。}

\begin{figure}[htbp]
\centering
\begin{tikzpicture}[>=stealth, scale=0.8]
    % 原始权重
    \node[draw, rectangle, fill=blue!20, minimum width=2.5cm, minimum height=2.5cm] (W) at (0, 0) {$W$};
    \node[below, font=\small] at (0, -1.5) {$d \times d$};
    \node[below, font=\scriptsize, gray] at (0, -2) {（冻结）};
    
    \node at (2, 0) {$+$};
    
    % B矩阵
    \node[draw, rectangle, fill=green!30, minimum width=0.5cm, minimum height=2.5cm] (B) at (4, 0) {$B$};
    \node[below, font=\small] at (4, -1.5) {$d \times r$};
    
    \node at (5, 0) {$\times$};
    
    % A矩阵
    \node[draw, rectangle, fill=orange!30, minimum width=2.5cm, minimum height=0.5cm] (A) at (7, 0) {$A$};
    \node[below, font=\small] at (7, -0.6) {$r \times d$};
    
    \node at (9.5, 0) {$=$};
    
    % 新权重
    \node[draw, rectangle, fill=purple!20, minimum width=2.5cm, minimum height=2.5cm] (Wnew) at (12, 0) {$W'$};
    \node[below, font=\small] at (12, -1.5) {$d \times d$};
    
    % 标注
    \draw[<->, red, thick] (3.7, -1.25) -- (4.3, -1.25);
    \node[red, below, font=\scriptsize] at (4, -1.4) {很窄!};
\end{tikzpicture}
\caption{LoRA：冻结原权重$W$，只训练低秩矩阵$A$和$B$}
\end{figure}

% ------------------------------------------------------------------------------------
\subsection*{参数量与显存}
% ------------------------------------------------------------------------------------

\begin{example}[LoRA参数量计算]
\textbf{设置}：$W$的形状是$[4096, 4096]$，LoRA秩$r=8$

\textbf{计算}：
\begin{align}
    \text{原参数量} &= 4096 \times 4096 = 16,777,216 \approx 16.7\text{M} \\
    \text{LoRA参数} &= 4096 \times 8 + 8 \times 4096 = 65,536 \approx 65\text{K} \\
    \text{比例} &= \frac{65\text{K}}{16.7\text{M}} \approx 0.4\%
\end{align}
\end{example}

LoRA的显存优势：

\begin{itemize}
    \item 原模型\textbf{冻结}：不需要存梯度和优化器状态
    \item 只训练$A$和$B$：参数量小，相应的梯度和优化器状态也小
\end{itemize}

\begin{table}[htbp]
\centering
\renewcommand{\arraystretch}{1.3}
\begin{tabular}{l|c|c}
\toprule
\textbf{方法} & \textbf{7B模型训练显存} & \textbf{单卡可行性} \\
\midrule
全量微调 & $\approx$100 GB & A100-80GB $\times$ \\
LoRA ($r=8$) & $\approx$18 GB & RTX 3090 \checkmark{} \\
QLoRA (4-bit + LoRA) & $\approx$6 GB & RTX 3080 \checkmark{} \\
\bottomrule
\end{tabular}
\caption{不同微调方法的显存需求对比}
\end{table}

% ------------------------------------------------------------------------------------
\subsection*{LoRA的关键细节}
% ------------------------------------------------------------------------------------

\paragraph{初始化}

\begin{itemize}
    \item $A$：高斯随机初始化
    \item $B$：\textbf{全零}初始化
\end{itemize}

为什么$B$初始化为0？因为这样初始时$BA = 0$，网络输出与预训练模型完全相同。这保证了微调从一个好的起点开始。

\paragraph{推理时合并权重}

训练完成后，可以把LoRA权重合并到原权重中：

\begin{equation}
    W' = W + BA
\end{equation}

合并后，推理时就是一个普通的线性层，\textbf{没有额外开销}。

\begin{lstlisting}[language=Python, caption=LoRA推理时合并权重]
# 训练完成后
merged_weight = base_weight + lora_B @ lora_A

# 推理时直接用merged_weight，速度与原模型相同
\end{lstlisting}

% ====================================================================================
\section*{高效计算：FlashAttention与MoE}
\addcontentsline{toc}{section}{高效计算：FlashAttention与MoE}
% ====================================================================================

\sidenoteplain{本节介绍两个现代大模型的关键技术：FlashAttention加速注意力计算，MoE降低推理成本。}

% ------------------------------------------------------------------------------------
\subsection*{FlashAttention：IO感知的注意力计算}
% ------------------------------------------------------------------------------------

\sidenoteplain{\textbf{FlashAttention}：2022年，斯坦福的Tri Dao在博士期间提出，通过IO感知算法将长序列注意力速度提升2-4倍，并大幅降低显存。该工作发表于NeurIPS 2022，Dao随后创办Together AI公司，FlashAttention已成为几乎所有大模型训练的标配。}

标准注意力计算有一个严重问题：需要存储完整的$N \times N$注意力矩阵。

\begin{keyformula}[title=标准Attention的显存瓶颈]
对于序列长度$N$，标准Attention需要：
\begin{itemize}
    \item 计算：$O(N^2 \cdot d)$ 次浮点运算
    \item 存储：$O(N^2)$ 的注意力矩阵（用于反向传播）
\end{itemize}

当$N = 100K$时，注意力矩阵需要 $100K \times 100K \times 2B = 20$ GB 显存！
\end{keyformula}

FlashAttention的核心思想：\textbf{分块计算 + 重算而非存储}。

\paragraph{算法思路}

\begin{enumerate}
    \item 将$Q, K, V$分成小块（block），每块能装进GPU的高速缓存（SRAM）
    \item 一次只计算一个块的注意力，立即与$V$相乘得到部分输出
    \item 使用在线softmax技巧，边算边累加，无需存储完整的注意力矩阵
    \item 反向传播时，重新计算注意力矩阵（recomputation），用计算换显存
\end{enumerate}

\begin{figure}[htbp]
\centering
\begin{tikzpicture}[scale=0.8]
    % 左边：标准Attention
    \node[font=\bfseries] at (0, 3.5) {标准Attention};

    % Q, K矩阵
    \draw[fill=blue!20] (0, 0) rectangle (1, 2.5);
    \node at (0.5, 2.8) {\small Q};
    \draw[fill=green!20] (1.5, 0) rectangle (4, 0.8);
    \node at (2.75, 1.1) {\small $K^T$};

    % 大的Attention矩阵
    \draw[fill=red!30] (0, -3) rectangle (2.5, -0.5);
    \node at (1.25, -1.75) {\small Attn};
    \node[font=\scriptsize] at (1.25, -2.3) {$N \times N$};
    \node[font=\scriptsize, red] at (1.25, -3.4) {需要存储！};

    % 箭头
    \draw[->, thick] (2, 0.4) -- (1.25, -0.3);

    % 右边：FlashAttention
    \node[font=\bfseries] at (7, 3.5) {FlashAttention};

    % 分块的Q, K
    \draw[fill=blue!40] (6, 1.8) rectangle (6.8, 2.5);
    \draw[fill=blue!20] (6, 1) rectangle (6.8, 1.7);
    \draw[fill=blue!20] (6, 0.2) rectangle (6.8, 0.9);
    \node at (6.4, 2.8) {\small Q块};

    \draw[fill=green!40] (7.2, 2.2) rectangle (9, 2.5);
    \draw[fill=green!20] (7.2, 1.9) rectangle (9, 2.2);
    \draw[fill=green!20] (7.2, 1.6) rectangle (9, 1.9);
    \node at (8.1, 2.8) {\small $K^T$块};

    % 小的块矩阵
    \draw[fill=orange!40] (6.5, -1.5) rectangle (7.3, -0.7);
    \node[font=\scriptsize] at (6.9, -1.1) {小块};
    \node[font=\scriptsize, green!60!black] at (6.9, -2) {在SRAM中};

    % 输出
    \draw[fill=purple!30] (8, -1.5) rectangle (8.6, -0.7);
    \node[font=\scriptsize] at (8.3, -1.1) {累加};

    % 箭头
    \draw[->, thick] (7.5, 1.6) -- (6.9, -0.5);
    \draw[->, thick] (7.3, -1.1) -- (7.9, -1.1);
\end{tikzpicture}
\caption{FlashAttention通过分块计算避免存储完整注意力矩阵}
\end{figure}

\begin{keyidea}[FlashAttention的核心洞察]
\textbf{GPU的瓶颈不是计算，而是数据搬运}（memory-bound）。

\begin{itemize}
    \item GPU算力：H100有990 TFLOPS
    \item 显存带宽：H100只有3.35 TB/s
    \item 计算/带宽比 = 296 FLOP/Byte，意味着每搬1字节要做296次运算才能喂饱GPU
\end{itemize}

FlashAttention把数据留在高速SRAM中，减少HBM读写，从而同时加速计算和节省显存。
\end{keyidea}
\clarify{"计算 / 带宽比"叫算术强度。Decode 阶段每读 2 字节权重只做 2 次运算（约 1 FLOP/Byte），远低于 GPU 的平衡点，所以 decode 是带宽受限——带宽约束是紧的，算力约束是松的。}

\paragraph{使用方法}

\begin{lstlisting}[language=Python, caption=PyTorch 2.0+ 内置FlashAttention]
import torch.nn.functional as F

# PyTorch 2.0+ 自动使用FlashAttention
output = F.scaled_dot_product_attention(
    query, key, value,
    is_causal=True  # 因果掩码
)

# 手动指定后端（PyTorch 2.3+ 请改用 torch.nn.attention.sdpa_kernel）
with torch.backends.cuda.sdp_kernel(
    enable_flash=True,
    enable_math=False,
    enable_mem_efficient=False
):
    output = F.scaled_dot_product_attention(query, key, value)
\end{lstlisting}

\begin{physicalmeaning}
\textbf{为什么FlashAttention能又快又省？}

回想一下，标准Attention需要：
\begin{enumerate}
    \item 计算 $n \times n$ 的注意力矩阵（$O(n^2)$ 空间）
    \item 把这个矩阵从SRAM $\to$ HBM $\to$ SRAM来回搬运
\end{enumerate}

FlashAttention的技巧：
\begin{itemize}
    \item 把Q分成小块，每次只处理一小块
    \item 在SRAM中算好softmax的增量更新，不存完整矩阵
    \item 结果：空间从 $O(n^2)$ 降到 $O(n)$，速度提升2-4倍
\end{itemize}

\textbf{对本科生的启示}：理解计算复杂度只是第一步，\textbf{理解数据搬运成本}才能真正优化程序。这就是为什么学习计算机体系结构（缓存、内存层次）很重要！
\end{physicalmeaning}

% ------------------------------------------------------------------------------------
\subsection*{Mixture of Experts (MoE)：稀疏激活}
% ------------------------------------------------------------------------------------

\sidenoteplain{\textbf{MoE的复兴}：MoE思想源自1991年Jacobs等人，但真正在大模型中大放异彩是Google的Shazeer(2017)在机器翻译中的应用。2023年底，Mistral AI发布的Mixtral 8x7B证明了开源MoE模型的巨大潜力。据传GPT-4也是MoE架构。}

MoE是另一种提升模型能力同时控制计算量的技术，被GPT-4、Mixtral等模型采用。

\begin{keyidea}[MoE的核心思想]
\textbf{模型参数很多，但每次只用一小部分}。

标准Transformer的FFN层对每个token都做相同的计算。MoE用多个"专家"（Expert）替代单个FFN，但每个token只路由到少数几个专家。
\end{keyidea}
\tryit{按激活参数 13B，用"计算量估算"一节的公式算 Mixtral 8x7B 每 token 的训练 FLOPs；再按 KV Cache 公式算它的缓存（32 层、GQA 8 个 KV 头、head\_dim = 128）。你会发现 MoE 省算力，不省显存。}

\paragraph{架构对比}

\begin{figure}[htbp]
\centering
\begin{tikzpicture}[scale=0.85]
    % 左边：标准FFN
    \node[font=\bfseries] at (1.5, 4) {标准FFN};

    \draw[fill=blue!20, rounded corners] (0, 2) rectangle (3, 3);
    \node at (1.5, 2.5) {FFN};

    \draw[fill=gray!20] (1.2, 0.5) rectangle (1.8, 1.5);
    \node[font=\scriptsize] at (1.5, 1) {token};

    \draw[->, thick] (1.5, 1.5) -- (1.5, 2);
    \draw[->, thick] (1.5, 3) -- (1.5, 3.5);

    \node[font=\scriptsize, gray] at (1.5, -0.3) {100\%参数激活};

    % 右边：MoE
    \node[font=\bfseries] at (8, 4) {MoE层};

    % 专家们
    \draw[fill=green!30, rounded corners] (5.5, 2) rectangle (6.5, 3);
    \node[font=\scriptsize] at (6, 2.5) {E1};

    \draw[fill=green!30, rounded corners] (6.8, 2) rectangle (7.8, 3);
    \node[font=\scriptsize] at (7.3, 2.5) {E2};

    \draw[fill=orange!40, rounded corners] (8.1, 2) rectangle (9.1, 3);
    \node[font=\scriptsize] at (8.6, 2.5) {E3};

    \draw[fill=green!30, rounded corners] (9.4, 2) rectangle (10.4, 3);
    \node[font=\scriptsize] at (9.9, 2.5) {E4};

    % Router
    \draw[fill=yellow!40, rounded corners] (7.3, 0.8) rectangle (8.7, 1.4);
    \node[font=\scriptsize] at (8, 1.1) {Router};

    % token
    \draw[fill=gray!20] (7.7, 0) rectangle (8.3, 0.5);
    \node[font=\scriptsize] at (8, 0.25) {token};

    % 箭头
    \draw[->, thick] (8, 0.5) -- (8, 0.8);
    \draw[->, thick, orange] (8, 1.4) -- (8.6, 2);
    \draw[->, thick, gray, dashed] (8, 1.4) -- (6, 2);
    \draw[->, thick, gray, dashed] (8, 1.4) -- (7.3, 2);
    \draw[->, thick, gray, dashed] (8, 1.4) -- (9.9, 2);

    \node[font=\scriptsize, gray] at (8, -0.5) {只激活Top-K专家};
    \node[font=\scriptsize, gray] at (8, -1) {(如K=2，激活25\%参数)};
\end{tikzpicture}
\caption{MoE用路由器选择性激活专家，减少计算量}
\end{figure}

\begin{keyformula}[title=MoE计算公式]
设有$E$个专家，每个token选择Top-K个：

\begin{equation}
    \text{Router}(x) = \text{softmax}(W_g \cdot x) \in \mathbb{R}^E
\end{equation}

\begin{equation}
    y = \sum_{i \in \text{TopK}} g_i \cdot \text{Expert}_i(x)
\end{equation}

其中$g_i$是门控权重，$\text{TopK}$是得分最高的K个专家。
\end{keyformula}

\paragraph{MoE的优势与挑战}

\begin{table}[htbp]
\centering
\renewcommand{\arraystretch}{1.3}
\begin{tabular}{l|l}
\toprule
\textbf{优势} & \textbf{挑战} \\
\midrule
参数量大但计算量可控 & 负载均衡问题（部分专家过载） \\
更强的模型容量 & 训练不稳定 \\
推理时只加载活跃专家 & 通信开销（分布式时） \\
\bottomrule
\end{tabular}
\end{table}

\begin{example}[Mixtral 8x7B的参数]
Mixtral 8x7B：
\begin{itemize}
    \item 总参数量：47B（8个7B专家 + 共享层）
    \item 激活参数量：约13B（每次只用2个专家）
    \item 计算量相当于13B模型，但能力接近大模型
\end{itemize}
\end{example}

\paragraph{负载均衡Loss}

为了避免所有token都路由到同一个专家，需要加入辅助损失：

\begin{equation}
    L_{\text{aux}} = \alpha \cdot E \cdot \sum_{i=1}^{E} f_i \cdot P_i
\end{equation}

其中$f_i$是分配给专家$i$的token比例，$P_i$是路由到专家$i$的平均概率。这个损失鼓励均匀分配。

% ====================================================================================
\section*{本附录小结}
\addcontentsline{toc}{section}{本附录小结}
% ====================================================================================

\begin{enumerate}
    \item \textbf{三种硬约束}：显存 $\le$ 显卡容量、训练时间 $\le$ 机时预算、数据搬运 $\le$ 带宽。目标（损失）没变，变的是约束。
    \item \textbf{显存估算}：推理 $\approx 2N$（FP16）+ KV Cache；全量微调 $\approx 12N$ + 激活值；LoRA 只训练低秩增量，显存降到 $2N$ 左右。
    \item \textbf{时间估算}：训练 FLOPs $\approx 6ND$，除以（GPU 数 $\times$ 算力 $\times$ MFU）。
    \item \textbf{推理吞吐}：decode 阶段带宽受限，tokens/s $\approx$ 带宽 / 模型大小。
    \item \textbf{所有"技巧"都是资源交换}：检查点与 FlashAttention 用算力买显存，KV Cache 用显存买算力，LoRA 用秩约束买显存，MoE 用参数总量买每 token 算力。
\end{enumerate}

\subsection*{核心公式速查}

\begin{table}[htbp]
\centering
\small
\renewcommand{\arraystretch}{1.5}
\begin{tabular}{l|l}
\toprule
\textbf{场景} & \textbf{公式} \\
\midrule
模型显存（FP16） & $M = 2N$ GB（$N$为参数量B） \\
全量微调显存 & $M \approx 12N$ GB（Adam优化器） \\
KV Cache & $M = 2 \times L \times S \times H \times D_{type}$ \\
\midrule
训练计算量 & FLOPs $\approx 6N$（每 token），总量 $6ND$ \\
训练时间 & $T = \frac{\text{总FLOPs}}{\text{GPU数} \times \text{算力} \times \text{MFU}}$ \\
推理吞吐量（Decode） & tokens/s $\approx \frac{\text{带宽}}{\text{模型大小}}$ \\
\midrule
Attention & $\text{softmax}\left(\frac{QK^\top}{\sqrt{d_k}}\right)V$ \\
LoRA参数量 & $(d_{in} + d_{out}) \times r$ \\
\bottomrule
\end{tabular}
\caption{本附录核心公式}
\end{table}

\subsection*{实战检查清单}

在开始实验前，问自己这几个问题：

\begin{enumerate}
    \item \textbf{显存够吗？}
    \begin{itemize}
        \item 推理：模型大小 + KV Cache
        \item 训练：$\times 6$（全量微调）或$\times 1.2$（LoRA）
    \end{itemize}
    
    \item \textbf{要跑多久？}
    \begin{itemize}
        \item 估算总FLOPs
        \item 除以（GPU数 $\times$ 算力 $\times$ 利用率）
    \end{itemize}
    
    \item \textbf{精度选什么？}
    \begin{itemize}
        \item 训练：BF16（如果支持）> FP16 + loss scaling
        \item 推理：FP16，显存紧张用INT8/INT4
    \end{itemize}
    
    \item \textbf{Batch Size怎么设？}
    \begin{itemize}
        \item 尽量大（提高GPU利用率）
        \item 但不能超过显存限制
        \item 可以用梯度累积模拟大Batch
    \end{itemize}
\end{enumerate}

\begin{physicalmeaning}[title=$\lambda$ 在本附录的意义：资源的影子价格]
本附录的 $\lambda$ 不再是梯度，而是\textbf{每种资源的边际价值}：
\[
\lambda_{\text{mem}} = -\frac{\partial\,\text{最优损失}}{\partial M_{\text{GPU}}}, \qquad
\lambda_{\text{time}} = -\frac{\partial\,\text{最优损失}}{\partial T_{\text{预算}}}, \qquad
\lambda_{\text{bw}} = -\frac{\partial\,\text{最优损失}}{\partial\,\text{带宽}}.
\]
\begin{itemize}
    \item \textbf{OOM} 就是显存约束绷紧、$\lambda_{\text{mem}}$ 很大：这时最值钱的不是更快的卡，而是更大的显存或更省显存的方法（量化、LoRA、检查点）。
    \item \textbf{梯度检查点与 FlashAttention 的重算}是在按汇率 $\lambda_{\text{mem}}/\lambda_{\text{time}}$ 用算力买显存；显存充裕（$\lambda_{\text{mem}} \approx 0$）时，重算只是白白浪费算力。
    \item \textbf{Decode 阶段带宽受限}意味着 $\lambda_{\text{bw}} > 0$ 而 $\lambda_{\text{time}} = 0$：再堆算力也没用，该做的是量化权重、量化 KV Cache 或加大 batch 提高算术强度。
    \item \textbf{MFU 只有 50\%}说明算力约束是松的——真正的紧约束在数据加载、通信或显存带宽上。
\end{itemize}
这和第1章的影子价格、第11章隐状态的影子价格 $\lambda_t$ 是同一个概念，只是这次的"约束"是你手里的显卡。第14章起实验规模都不大，但先算一遍显存和时间，是每次开跑前的习惯。
\end{physicalmeaning}

% ====================================================================================
\section*{习题}
\addcontentsline{toc}{section}{习题}
% ====================================================================================

\paragraph{计算题}

\begin{enumerate}
    \item 你有一张RTX 4090（24GB），想跑LLaMA-13B推理。
    \begin{enumerate}
        \item FP16能跑吗？
        \item 如果不能，INT8呢？INT4呢？
        \item 假设用INT4，序列长度4K，Batch=4时KV Cache占多少显存？能跑吗？
    \end{enumerate}
    
    \item 实验室有8张A100-40GB。
    \begin{enumerate}
        \item 能否全量微调LLaMA-7B？如果不能，需要什么技术？
        \item 如果用LoRA，单卡能跑吗？
    \end{enumerate}
    
    \item 计算以下配置的KV Cache显存：
    \begin{itemize}
        \item 模型：48层，hidden\_size = 8192
        \item 序列长度：32K
        \item Batch Size：16
        \item 精度：FP16
    \end{itemize}
\end{enumerate}

\begin{enumerate}
    \setcounter{enumi}{3}
    
    \item 用50K条数据（平均长度1024 tokens）对7B模型做LoRA微调，训练5个epoch。
    \begin{enumerate}
        \item 估算总FLOPs
        \item 用单张RTX 4090（83 TFLOPS，MFU=45\%）需要多久？
        \item 用4张A100（312 TFLOPS，MFU=50\%）需要多久？
    \end{enumerate}
    
    \item LLaMA-7B在A100-80GB上的推理吞吐量约为90 tokens/s。你需要生成1000条回复，每条平均500 tokens。需要多久？
\end{enumerate}

\paragraph{编程题}

\begin{enumerate}
    \setcounter{enumi}{5}
    
    \item 实现 FFN（含 GELU）和 LayerNorm，与本附录的 \texttt{MultiHeadAttention} 组装成完整的 Pre-Norm \texttt{TransformerBlock}，用 \texttt{torch.autograd.gradcheck} 验证梯度。
    
    \item 实现因果掩码的生成函数\texttt{create\_causal\_mask(seq\_len)}。
    
    \item 给本附录的 \texttt{MultiHeadAttention} 加上 KV Cache，测量序列长度 1K / 4K / 16K 时逐 token 生成的时间与显存，验证"用显存买算力"。
\end{enumerate}

\paragraph{思考题}

\begin{enumerate}
    \setcounter{enumi}{8}
    
    \item 为什么长文本推理（如100K context）特别吃显存？有哪些缓解方法？
    
    \item LoRA的$B$矩阵为什么初始化为0？如果用高斯初始化会怎样？
    
    \item Flash Attention为什么能同时加速计算和节省显存？（提示：分块计算、IO感知）
\end{enumerate}


% ====================================================================================
% 本附录原有的"扩展阅读：从零上手——服务器与环境配置"已整体移至附录B（实验环境速成），
% 以免与 Transformer 的数学主线混在一起。
% ====================================================================================
\paragraph{关于实验环境}
本附录的代码需要在带 GPU 的服务器上运行。第一次连接服务器、配置 conda 环境、用 tmux 挂后台任务的同学，请先读附录B"实验环境速成——服务器与环境配置"，那里从 SSH 登录讲到提交训练任务，随用随查。


\section*{深度学习的计算基础设施}
\addcontentsline{toc}{section}{计算基础设施}

当你写下 \texttt{y = torch.matmul(A, B)} 时，背后发生了什么？为什么 GPU 
比 CPU 快 100 倍？为了理解这一点，我们先从一个经典的小故事开始。

\subsection*{一个被认为“毫无意义”的实验}

\sidenoteplain{\textbf{Geoffrey Hinton}(1947--)：``深度学习教父''，2018年图灵奖得主。他在多伦多大学坚持神经网络研究三十年，在AI冬天时几乎是孤军奋战。2012年AlexNet的成功证明了他的坚持。2023年他辞去Google职位，公开警告AI的潜在危险。}

2009 年，Montreal 的一间小实验室里，博士生 Ranzato、
助手 Alex Krizhevsky，以及他们的导师 Hinton 正在做一个
很多人看来``一点用都没有''的实验——

\begin{quote}
用当时游戏玩家的显卡（NVIDIA GTX 280）来训练一个
卷积神经网络。
\end{quote}

在那个年代，深度学习还远没有今天的气势。主流观点认为：

\begin{itemize}
    \item 深度网络太大，训练太慢；
    \item GPU 只是玩游戏的，没有科研价值；
    \item 机器学习的未来应该是从数学上给模型加约束，而不是暴力堆算力。
\end{itemize}

但 Krizhevsky 注意到了一件别人忽略的事：  
游戏显卡的核心结构（SIMD + 大量并行浮点单元）和卷积操作的数学结构是完美匹配的。

于是他们写下了第一版用 CUDA 加速卷积的代码，
发现速度比 CPU \textbf{快了几十倍}。Hinton 后来回忆说：

\begin{quote}
``第一次看到 GPU 把卷积层跑得这么快时，我们觉得好像未来突然被打开了。''
\end{quote}

三年后，也就是 2012 年，Krizhevsky 基于这些经验写出了 \texttt{cuda-convnet}，并用一台双 GTX 580 显卡训练出了历史性的模型——

\begin{center}
\textbf{AlexNet}
\end{center}

它把 ImageNet 错误率降低了 10 个百分点，直接点燃了整个深度学习时代。

从此之后，一个事实被彻底确认：

\begin{quote}
\textbf{深度学习的突破，离不开硬件架构与数学结构的完美匹配。}
\end{quote}

而你今天执行的 \texttt{torch.matmul}、\texttt{conv2d} 等操作，
本质上都是在 GPU 上调用几十年积累下来的线性代数内核（cuBLAS、cuDNN），
这些内核又精确地利用了 GPU 的大规模并行结构。

因此，表面上一行 Python：

\begin{equation}
y = \text{torch.matmul}(A, B)
\end{equation}

背后是一整套工业级的软硬件生态：计算图编译器、Tensor Core、
warp 调度、内存对齐、缓存优化、并行线程块调度……

理解这一点后，PyTorch 为什么快，就不再是魔法，而是工程与数学共同的必然结果(这一节以英伟达NVIDIA相关生态举例，华为等硬件厂商也有类似的底层架构逻辑，机器码和线程管理模式不尽相同)。

\subsection*{GPU 编程的层次结构}

深度学习的计算栈是这样的：

\begin{figure}[htbp]
\centering
\begin{tikzpicture}[
    box/.style={draw, rounded corners, minimum width=4cm, minimum height=0.9cm, font=\small, align=center},
    arrow/.style={->, thick, >=stealth}
]
    \node[box, fill=blue!20] (python) at (0, 5) {Python 代码\\{\ttfamily y = model(x)}};
    \node[box, fill=blue!30] (pytorch) at (0, 3.8) {PyTorch / JAX / TensorFlow\\自动微分、张量操作};
    \node[box, fill=green!20] (cuda) at (0, 2.6) {CUDA / cuDNN / cuBLAS\\GPU 算子库};
    \node[box, fill=yellow!20] (driver) at (0, 1.4) {GPU 驱动\\内存管理、任务调度};
    \node[box, fill=red!20] (hw) at (0, 0.2) {GPU 硬件\\CUDA Core / Tensor Core};
    
    \draw[arrow] (python) -- (pytorch);
    \draw[arrow] (pytorch) -- (cuda);
    \draw[arrow] (cuda) -- (driver);
    \draw[arrow] (driver) -- (hw);
    
    % 右侧标注
    \node[right, font=\small, text=gray] at (2.3, 5) {你写的};
    \node[right, font=\small, text=gray] at (2.3, 3.8) {框架提供};
    \node[right, font=\small, text=gray] at (2.3, 2.6) {NVIDIA 提供};
    \node[right, font=\small, text=gray] at (2.3, 1.4) {系统层};
    \node[right, font=\small, text=gray] at (2.3, 0.2) {硬件层};
\end{tikzpicture}
\caption{深度学习计算栈：从 Python 到硬件}
\end{figure}

\subsection*{CUDA：GPU 编程的基础}

\textbf{CUDA}（Compute Unified Device Architecture）是 NVIDIA 的 GPU 编程平台。

\paragraph{核心概念}

\begin{itemize}
    \item \textbf{Thread}：最小执行单元，执行相同的代码（SIMT: Single Instruction Multiple Threads）
    \item \textbf{Block}：一组线程，共享高速的 Shared Memory
    \item \textbf{Grid}：所有 Block 的集合，完成一个计算任务
\end{itemize}

\begin{figure}[htbp]
\centering
\begin{tikzpicture}[scale=0.8]
    % Grid
    \draw[thick, blue] (0, 0) rectangle (8, 4);
    \node[blue, font=\bfseries] at (4, 4.4) {Grid（整个任务）};
    
    % Blocks
    \foreach \i in {0, 2.7, 5.4} {
        \draw[thick, green!60!black, fill=green!10] (\i+0.2, 0.2) rectangle (\i+2.5, 3.8);
    }
    \node[green!60!black, font=\small] at (1.35, 3.5) {Block 0};
    \node[green!60!black, font=\small] at (4.05, 3.5) {Block 1};
    \node[green!60!black, font=\small] at (6.75, 3.5) {Block 2};
    
    % Threads in first block
    \foreach \i in {0,...,3} {
        \foreach \j in {0,...,2} {
            \draw[fill=red!30] (0.4+\i*0.5, 0.5+\j*0.9) rectangle ++(0.4, 0.7);
        }
    }
    \node[font=\scriptsize] at (1.35, 0.3) {Threads};
    
    % 标注
    \node[right, font=\small, align=left] at (8.5, 3) {每个 Block：\\共享 Shared Memory\\同步执行};
    \node[right, font=\small, align=left] at (8.5, 1) {每个 Thread：\\独立寄存器\\执行相同代码};
\end{tikzpicture}
\caption{CUDA 的线程层次：Grid → Block → Thread}
\end{figure}

\paragraph{一个简单的 CUDA 核函数}

\begin{lstlisting}[style=teaching, language=C]
// 向量加法的 CUDA 实现
__global__ void vector_add(float *a, float *b, float *c, int n) {
    // 计算当前线程负责的元素索引
    int idx = blockIdx.x * blockDim.x + threadIdx.x;
    
    if (idx < n) {
        c[idx] = a[idx] + b[idx];  // 每个线程只算一个元素！
    }
}

// 调用：启动 N 个线程并行计算
vector_add<<<num_blocks, threads_per_block>>>(a, b, c, N);
\end{lstlisting}

\textbf{关键洞察}：传统 CPU 代码用 \texttt{for} 循环遍历数组，CUDA 代码让每个线程只处理一个元素，成千上万个线程\textbf{同时}执行！

\subsection*{内存层次与性能优化}

GPU 有多级内存，访问速度差异巨大：

\begin{center}
\renewcommand{\arraystretch}{1.4}
\begin{tabular}{l|r|r|l}
\toprule
\textbf{内存类型} & \textbf{大小} & \textbf{带宽} & \textbf{延迟} \\
\midrule
寄存器（Register） & $\sim$256 KB/SM & -- & 1 cycle \\
共享内存（Shared Memory） & $\sim$100 KB/SM & $\sim$19 TB/s & $\sim$20 cycles \\
L2 缓存 & $\sim$40 MB & $\sim$6 TB/s & $\sim$200 cycles \\
全局内存（Global/HBM） & 24--80 GB & $\sim$2 TB/s & $\sim$400 cycles \\
\bottomrule
\end{tabular}
\end{center}

\sidenoteplain{SM = Streaming Multiprocessor，GPU 的计算单元。一个 GPU 有几十到上百个 SM。}

\textbf{性能优化的核心}：尽量让数据停留在快速的内存层次中！

\begin{itemize}
    \item \textbf{访存合并}（Memory Coalescing）：相邻线程访问相邻内存，一次传输搞定
    \item \textbf{Shared Memory}：Block 内的线程共享数据，避免重复读取全局内存
    \item \textbf{寄存器复用}：计算密集型操作尽量用寄存器存中间结果
\end{itemize}

\subsection*{张量的内存布局：Stride 的秘密}

当你创建一个 PyTorch 张量时，它在内存中是怎么存储的？

\begin{lstlisting}[style=teaching, language=Python]
import torch

A = torch.tensor([[1, 2, 3],
                  [4, 5, 6]])  # shape: (2, 3)

print(A.storage())  # 底层存储：[1, 2, 3, 4, 5, 6]（一维数组！）
print(A.stride())   # 步长：(3, 1)
\end{lstlisting}

\textbf{Stride}（步长）告诉我们：沿每个维度移动一步，需要在底层数组中跳过多少元素。

\begin{figure}[htbp]
\centering
\begin{tikzpicture}[scale=0.9]
    % 底层一维数组
    \foreach \i/\v in {0/1, 1/2, 2/3, 3/4, 4/5, 5/6} {
        \draw[fill=blue!20] (\i*1.2, 0) rectangle ++(1, 0.8);
        \node at (\i*1.2+0.5, 0.4) {\v};
        \node[font=\scriptsize, gray] at (\i*1.2+0.5, -0.3) {\i};
    }
    \node[font=\small] at (3.5, -0.8) {底层存储（连续内存）};
    
    % 逻辑视图
    \node at (-1.5, 2.5) {逻辑视图：};
    \draw[fill=green!20] (0, 2) rectangle ++(1, 0.8); \node at (0.5, 2.4) {1};
    \draw[fill=green!20] (1.2, 2) rectangle ++(1, 0.8); \node at (1.7, 2.4) {2};
    \draw[fill=green!20] (2.4, 2) rectangle ++(1, 0.8); \node at (2.9, 2.4) {3};
    \draw[fill=green!20] (0, 3) rectangle ++(1, 0.8); \node at (0.5, 3.4) {4};
    \draw[fill=green!20] (1.2, 3) rectangle ++(1, 0.8); \node at (1.7, 3.4) {5};
    \draw[fill=green!20] (2.4, 3) rectangle ++(1, 0.8); \node at (2.9, 3.4) {6};
    
    % stride 标注
    \draw[<->, thick, red] (0.5, 1.8) -- (0.5, 1.2) -- (3.5, 1.2) -- (3.5, 0.9);
    \node[red, font=\small] at (2, 1.5) {stride[0]=3};
    
    \draw[<->, thick, orange] (4.5, 2.4) -- (5.5, 2.4);
    \node[orange, font=\small] at (5, 2.8) {stride[1]=1};
\end{tikzpicture}
\caption{张量的 Stride：逻辑多维数组如何映射到一维内存}
\end{figure}

\paragraph{为什么 Stride 重要？}

\textbf{转置}不需要复制数据，只需要交换 stride！

\begin{lstlisting}[style=teaching, language=Python]
A = torch.tensor([[1, 2, 3],
                  [4, 5, 6]])
print(A.stride())       # (3, 1) - 行优先

B = A.T                 # 转置
print(B.stride())       # (1, 3) - 列优先
print(B.storage())      # 还是 [1, 2, 3, 4, 5, 6]，没有复制！

# 但是 B 不是"连续"的
print(B.is_contiguous())  # False
\end{lstlisting}

\begin{insightbox}[title=Contiguous 的含义]
一个张量是 \textbf{contiguous} 的，当且仅当它的内存布局与从头创建时一致（行优先，stride 递减）。

很多 CUDA 算子要求输入是 contiguous 的，否则需要先调用 \texttt{.contiguous()} 复制一份数据。这是常见的性能陷阱！
\end{insightbox}

\subsection*{cuBLAS 和 cuDNN：优化到极致的算子库}

自己写 CUDA 核函数很难达到最优性能。NVIDIA 提供了高度优化的库：

\begin{itemize}
    \item \textbf{cuBLAS}：基础线性代数（矩阵乘法、向量运算）
    \item \textbf{cuDNN}：深度学习原语（卷积、BatchNorm、RNN）
    \item \textbf{cuFFT}：快速傅里叶变换
    \item \textbf{cuSPARSE}：稀疏矩阵运算
\end{itemize}

\textbf{矩阵乘法有多复杂？}
\sidenoteplain{一个优化良好的矩阵乘法实现可能有上万行代码，针对不同尺寸有几十个特化版本。}

看似简单的 $C = AB$，cuBLAS 的实现考虑了：
\begin{itemize}
    \item 矩阵分块（Tiling）：适配 Shared Memory 大小
    \item 数据预取（Prefetching）：隐藏内存延迟
    \item 向量化访存：一次读取 128 位
    \item Tensor Core 加速：专用硬件做 $4 \times 4$ 矩阵乘
    \item 不同矩阵尺寸的特化实现
\end{itemize}



\subsection*{Triton：让 GPU 编程更简单}

\textbf{Triton} 是由 OpenAI 推出的新一代 GPU 编程语言，旨在让研究者能够\textbf{以接近 Python 的方式编写高性能的自定义算子}，而无需深度掌握 CUDA 的底层细节。它在矩阵乘法、注意力等核心算子上曾展现出与手写 CUDA 相当甚至更优的性能。由于长期维护成本高、生态兼容性复杂，Triton 已于 2025 年 11 月正式停止维护，但其思想（如自动优化、块级并行抽象）仍然对后续深度学习编译器的设计产生了重要影响。

\begin{lstlisting}[style=teaching, language=Python]
import triton
import triton.language as tl

@triton.jit
def vector_add_kernel(x_ptr, y_ptr, out_ptr, n, BLOCK_SIZE: tl.constexpr):
    # 每个程序实例处理 BLOCK_SIZE 个元素
    pid = tl.program_id(0)
    offsets = pid * BLOCK_SIZE + tl.arange(0, BLOCK_SIZE)
    mask = offsets < n
    
    # 向量化加载、计算、存储
    x = tl.load(x_ptr + offsets, mask=mask)
    y = tl.load(y_ptr + offsets, mask=mask)
    tl.store(out_ptr + offsets, x + y, mask=mask)
\end{lstlisting}

\textbf{Triton vs CUDA}：

\begin{center}
\renewcommand{\arraystretch}{1.4}
\begin{tabular}{l|l|l}
\toprule
\textbf{方面} & \textbf{CUDA} & \textbf{Triton} \\
\midrule
语言 & C++ 扩展 & Python 风格 \\
抽象层次 & 线程级 & Block 级 \\
内存管理 & 手动 & 自动 Tiling \\
性能调优 & 繁琐 & 编译器自动优化 \\
学习曲线 & 陡峭 & 相对平缓 \\
\bottomrule
\end{tabular}
\end{center}

PyTorch 2.0 的 \texttt{torch.compile} 就使用 Triton 生成高效的融合算子。

\subsection*{算子融合：减少内存访问}

深度学习中一个关键优化是\textbf{算子融合}（Kernel Fusion）。

考虑 $y = \text{ReLU}(Wx + b)$，朴素实现：

\begin{lstlisting}[style=teaching, language=Python]
# 朴素实现：3 次 kernel launch，3 次读写全局内存
z1 = torch.matmul(W, x)   # 写 z1 到 HBM
z2 = z1 + b               # 读 z1，写 z2 到 HBM  
y = torch.relu(z2)        # 读 z2，写 y 到 HBM
\end{lstlisting}

\textbf{问题}：每个操作都要把中间结果写回慢速的全局内存（HBM），再读出来。

\textbf{融合后}：一个 kernel 完成所有操作，中间结果只在寄存器/Shared Memory 中：

\begin{lstlisting}[style=teaching, language=Python]
# 融合实现：1 次 kernel launch，中间结果不写回 HBM
@triton.jit
def fused_linear_relu(x, W, b, y, ...):
    z1 = dot(W, x)        # 结果在寄存器
    z2 = z1 + b           # 结果在寄存器
    y = max(z2, 0)        # 直接写最终结果
\end{lstlisting}

\begin{figure}[htbp]
\centering
\begin{tikzpicture}[scale=0.8]
    % 朴素版本
    \begin{scope}
        \node[font=\bfseries] at (2, 4) {朴素实现};
        \draw[fill=green!20] (0, 3) rectangle (1.5, 3.5); \node[font=\small] at (0.75, 3.25) {matmul};
        \draw[fill=green!20] (2, 3) rectangle (3.5, 3.5); \node[font=\small] at (2.75, 3.25) {add};
        \draw[fill=green!20] (4, 3) rectangle (5.5, 3.5); \node[font=\small] at (4.75, 3.25) {relu};
        
        \draw[fill=red!20] (0, 1.5) rectangle (5.5, 2.2);
        \node[font=\small] at (2.75, 1.85) {HBM（全局内存）};
        
        \draw[->, thick] (0.75, 3) -- (0.75, 2.2);
        \draw[->, thick] (0.75, 2.2) -- (2.75, 2.2) -- (2.75, 3);
        \draw[->, thick] (2.75, 3) -- (2.75, 2.2);
        \draw[->, thick] (2.75, 2.2) -- (4.75, 2.2) -- (4.75, 3);
        
        \node[font=\small, red] at (2.75, 0.8) {6 次 HBM 访问};
    \end{scope}
    
    % 融合版本
    \begin{scope}[xshift=8cm]
        \node[font=\bfseries] at (2, 4) {融合实现};
        \draw[fill=blue!30] (0.5, 2.8) rectangle (4, 3.7);
        \node[font=\small] at (2.25, 3.25) {fused\_linear\_relu};
        
        \draw[fill=yellow!30] (1, 2) rectangle (3.5, 2.5);
        \node[font=\small] at (2.25, 2.25) {寄存器};
        
        \draw[fill=red!20] (0, 0.8) rectangle (4.5, 1.3);
        \node[font=\small] at (2.25, 1.05) {HBM};
        
        \draw[->, thick, dashed] (2.25, 2.8) -- (2.25, 2.5);
        \draw[->, thick] (2.25, 2) -- (2.25, 1.3);
        
        \node[font=\small, blue] at (2.25, 0.3) {1 次 HBM 写入};
    \end{scope}
\end{tikzpicture}
\caption{算子融合减少内存访问：6 次 $\to$ 1 次}
\end{figure}

FlashAttention 就是通过融合 QKV 计算、Softmax、输出投影，避免了存储巨大的注意力矩阵。

\subsection*{深度学习软件生态}

\begin{figure}[htbp]
\centering
\begin{tikzpicture}[
    box/.style={draw, rounded corners, minimum height=0.8cm, font=\small, align=center},
    scale=0.85, transform shape
]
    % 应用层
    \node[box, fill=purple!20, minimum width=8cm] at (4, 5) {应用：ChatGPT, Stable Diffusion, AlphaFold};
    
    % 框架层
    \node[box, fill=blue!20, minimum width=2.4cm] at (1.2, 4) {PyTorch};
    \node[box, fill=blue!20, minimum width=2.4cm] at (4, 4) {JAX};
    \node[box, fill=blue!20, minimum width=2.4cm] at (6.8, 4) {TensorFlow};
    
    % 编译器层
    \node[box, fill=green!20, minimum width=2.4cm] at (1.2, 3) {TorchScript};
    \node[box, fill=green!20, minimum width=2.4cm] at (4, 3) {XLA};
    \node[box, fill=green!20, minimum width=2.4cm] at (6.8, 3) {Triton};
    
    % 运行时层
    \node[box, fill=yellow!20, minimum width=3.8cm] at (2, 2) {cuDNN / cuBLAS};
    \node[box, fill=yellow!20, minimum width=3.8cm] at (6, 2) {NCCL（多卡通信）};
    
    % 驱动层
    \node[box, fill=orange!20, minimum width=8cm] at (4, 1) {CUDA Runtime / Driver};
    
    % 硬件层
    \node[box, fill=red!20, minimum width=3.8cm] at (2, 0) {NVIDIA GPU};
    \node[box, fill=red!20, minimum width=3.8cm] at (6, 0) {AMD GPU / TPU};
\end{tikzpicture}
\caption{深度学习软件栈全景图}
\end{figure}

\begin{center}
\renewcommand{\arraystretch}{1.4}
\begin{tabular}{l|l}
\toprule
\textbf{组件} & \textbf{作用} \\
\midrule
PyTorch / JAX & 前端框架，提供自动微分、张量操作 \\
Triton / XLA & 编译器，生成优化的 GPU 代码 \\
cuDNN / cuBLAS & 算子库，提供高度优化的基础运算 \\
NCCL & 多 GPU 通信库，支持分布式训练 \\
CUDA & GPU 编程平台和运行时 \\
\bottomrule
\end{tabular}
\end{center}

\subsection*{为什么了解这些很重要？}

\begin{enumerate}
    \item \textbf{性能调优}：知道瓶颈在哪里（计算 vs 内存 vs 通信）
    \item \textbf{写高效代码}：避免不必要的 \texttt{.contiguous()}、合理使用 \texttt{torch.compile}
    \item \textbf{理解论文}：FlashAttention、Paged Attention 等优化的原理
    \item \textbf{职业发展}：GPU 编程是稀缺技能，薪资溢价高
\end{enumerate}

\begin{insightbox}[title=实用建议]
\begin{itemize}
    \item 入门：先用好 PyTorch 的 \texttt{torch.compile}，让编译器帮你优化
    \item 进阶：学习 Triton，能写自定义高效算子
    \item 深入：学习 CUDA，理解底层原理
    \item 关注：FlashAttention、vLLM 等开源项目的实现
\end{itemize}
\end{insightbox}


